\section{Introduction}

In type theory, universes allow manipulating types as terms.
To avoid Girard's paradox, Martin-Löf \cite{Martin-Lof-1972, MARTINLOF75} introduced a \textit{hierarchy} of universes: $\U_0, \U_1, \U_2, ...$.
Those can be assumed to be \textit{cumulative}, requiring proper inclusions between them.
Martin-Löf’s first presentations \cite{MARTINLOF75, MARTINLOF82} were formulated à la Russell, where a universe is seen as a type of types. Later, an alternative version à la Tarski \cite{MARTINLOF84}, where a universe is seen as a type of codes along with a decoding function, was also introduced. Usually, Russell universes are regarded as an informal notation of the latter, as \citet{Luo} suggested.

When presenting Type Theory in an algebraic fashion, Tarski universes are essential, at least if we do not want to introduce sort equations. Indeed, identifying the sort of terms of a universe with a sort of types forces strict set-theoretic equalities that fail in many semantic models. If people often assume that semantics result extend easily to the more informal syntax used in practice, this comes from misconceptions as they depend on subtle details of the presentation. To solve this issue, Voevodsky proposed the \textit{Initiality Conjecture project}, which aimed at developing rigorous metatheory for a large class of theories.
A first step at solving this goal was achieved by Streicher \cite{Streicher}, proving initiality for a fully annotated Calculus of Constructions, and then proving its equivalence with the unannotated formulation. 
Generalised Algebraic Theories (GATs) provide a framework to describe type theories and their models in a uniform way. However, they give rise to a fully explicit and annotated presentation of type theory, which is inconvenient to work in.
Therefore, it would be essential to extend Streicher's equivalence results to relate the Generalised Algebraic and Pure Type System presentations. This would justify rigorously the common practice of using informal unannotated syntax with Russell universes, and would legitimate interpreting such syntax in the usual models of type theory.

\noindent
\paragraph*{Related work}
Assaf \cite{assaf} showed the equivalence between Russell and Tarski presentations of universes for the Calculus of Constructions with cumulativity and $\beta$ as an external reduction. A proof extending this result to cumulativity with product covariance was recently published by Felicissimo and Winterhalter \cite{tiago-fscd}, using a technique of \citet{Streicher}. Unfortunately, because these approaches rely on confluence and view $\beta$ as an external reduction, their arguments do not apply to formulations of type theory with judgemental equality, and do not extend to $\eta$-equivalence. This justifies the study of a new proof technique, which applies directly to equational presentations of type theory with $\beta\eta$-conversion.

In the present paper, we will we will study the simplest form of cumulativity, differing from \citet{tiago-fscd} who retail product covariance. Indeed, structural sub-typing breaks canonicity when adding inductive types \cite{Luo}, and product contravariance does not hold in the set model. Also, it is essential to note that both product covariance and contravariance can be simulated with eta expansion in our setting. Yet, stronger forms of structural sub-typing would provide a more convenient manipulation of universes in practical systems \cite{subtyping}. The recent work of \citet{AdapTT} presents a general framework which would provide automatically the necessary equations to solve the issues mentioned above. In any case, we believe that our work is orthogonal to such considerations, and that our proof should extend to any system that handles structural sub-typing in a correct way.

\paragraph*{Contributions} In this paper we relate the Generalised Algebraic presentation of type theory to the implicit presentation inherited from Pure Type Systems, by a method due to \citet{Streicher}. We considered Martin-Löf's Type Theory with dependent products, cumulative universes and judgmental $\beta\eta$ conversion. Our proof is modular and adapts to various extensions of type theory, such as universe polymorphism \cite{polymorphism}.
Concretely, our equivalence result entails that any semantic model $\mathcal{M}$ of the GAT interprets derivations of the implicit, unannotated Russell theory.
This bridges the gap between the correction and semantics of explicit presentations, and the implicit theories actually implemented in proof systems.
Additionally, this relates modern semantic results to the literature where implicit syntax was the norm.

The main argument is as follows: there are forgetful maps from the implicit theories to the
implicit presentation, and we show that any implicit judgement can be lifted back in a unique way.
To prove this, we rely on well-known properties of the syntax : no-confusion and injectivity of type formers, which are sometimes referred to as \textit{definitional inversion} \cite{definitional-inv}.
We prove them for the generalised algebraic theory, and they will propagate along our equivalence results, differing from \citet{tiago-fscd} who need injectivity for each theories.
Those properties seem to be strongly related to conversion-checking \cite{meven}, and our equivalence result itself is to be related to the elaboration in proof assistants.

Our proof is elementary and provides an algorithm to translate from one presentation to the other. The definitional inversion properties are proven using the new domain-theoretic technique of \citet{definitional-inv}, making our proof adaptable to various theories which do not enjoy normalisation.
In fact, the proof is purely syntactic and can be done internal to a weak metatheory such as primitive recursive arithmetic, whereas normalisation relies on a stronger metatheory such as the constructive set theory of \citet{Coq19, aczel}.

We did a human-written formalisation of the equivalence between Tarski and Russell universes in Rocq \cite{Formalisation}, admitting definitional inversion properties. Throughout the paper, we provide links to this development, indicated by the Rocq logo \includegraphics[height=1em]{icon-rocq-blue-eps-converted-to.pdf}.
Along with this, we did a complete auto-formalisation in Agda of the whole paper using Claude \cite{auto-formalisation}. This was useful to ensure correctness and find some mistakes in our proofs. More precisely, we found an issue in the equivalence proof of \zcref{sectionT2T2'}, and we checked that our result extends to universe polymorphism (\zcref{extension}). We also auto-formalised the whole section on interpretation (\zcref{sec:semantics}) using Unit \cite{Coquand2026Unit}, allowing us to test our result on concrete models. We verified carefully that this faithfully matches our statements, and that no postulates were used.

\paragraph*{Structure of this paper}
In \zcref{GAT} we present the fully explicit Generalised Algebraic Theory and prove definitional inversion, before presenting the zoo of implicit presentations in \zcref{implicit-presentations}.
Then, we prove that all those syntaxes are equivalent in the following subsections: 
\vspace{0.5em}
\begin{itemize}
    \item[] \textbf{\S\S~\ref{sectionT0T1}.} We relate explicit and implicit substitutions. 
    \item[] \textbf{\S\S~\ref{sectionT1T2}.} We relate Tarski and Russell presentations for cumulative universes.
    \item[] \textbf{\S\S~\ref{sectionT2T2'}.} We relate annotated and unannotated syntaxes. 
\end{itemize}
\vspace{0.5em}
Finally, we extend initiality and give semantics to implicit syntax in \zcref{sec:semantics}.

